% Figure : bénéfice net moyen d'une politique par crise, selon le coût d'un levier
% (données : ip_timeseries.csv et base d'apprentissage du lot A ; calcul décrit en section 6.6)
\begin{tikzpicture}
\begin{axis}[
  width=0.93\textwidth, height=7.2cm,
  xmin=0, xmax=12000, ymin=-6, ymax=20,
  xlabel={Coût d'un levier engagé, exprimé en unités de marge ($c/m$)},
  ylabel={Bénéfice net moyen par crise (milliers d'unités de marge)},
  xtick={0,2000,4000,6000,8000,10000,12000},
  xticklabel style={/pgf/number format/.cd,1000 sep={\,}},
  scaled x ticks=false,
  tick label style={font=\scriptsize}, label style={font=\small},
  grid=major, grid style={black!8},
  legend style={font=\scriptsize, draw=none, fill=white, fill opacity=0.9, text opacity=1, at={(0.98,0.98)}, anchor=north east, cells={anchor=west}},
  every axis plot/.append style={line width=1.2pt}]
  \fill[polreac!7] (axis cs:0,-6) rectangle (axis cs:6628,20);
  \node[font=\scriptsize, text=polreac!80!black, anchor=south west, align=left] at (axis cs:150,-5.6) {zone où la politique\\réactive est rentable};
  \addplot[polnone] table[x=r, y expr=\thisrow{none}/1000] {figures/data/eco_benefice.dat};
  \addlegendentry{Sans réaction}
  \addplot[poltard, dotted] table[x=r, y expr=\thisrow{tardif}/1000] {figures/data/eco_benefice.dat};
  \addlegendentry{Tardif}
  \addplot[polprud, dashed] table[x=r, y expr=\thisrow{prudent}/1000] {figures/data/eco_benefice.dat};
  \addlegendentry{Prudent}
  \addplot[polreac] table[x=r, y expr=\thisrow{reactif}/1000] {figures/data/eco_benefice.dat};
  \addlegendentry{Réactif}
  \addplot[palmec, line width=1.6pt, densely dashed] table[x=r, y expr=\thisrow{posteriori}/1000] {figures/data/eco_benefice.dat};
  \addlegendentry{Meilleure politique choisie crise par crise}
  \draw[black!50, densely dotted] (axis cs:6628,-6) -- (axis cs:6628,20);
  \node[font=\scriptsize, fill=white, inner sep=1pt, anchor=west] at (axis cs:6750,-4.2) {seuil de rentabilité de la politique réactive : 6 630};
\end{axis}
\end{tikzpicture}
